% Options for packages loaded elsewhere
% Options for packages loaded elsewhere
\PassOptionsToPackage{unicode}{hyperref}
\PassOptionsToPackage{hyphens}{url}
\PassOptionsToPackage{dvipsnames,svgnames,x11names}{xcolor}
%
\documentclass[
  11pt,
]{article}
\usepackage{xcolor}
\usepackage[margin=1in]{geometry}
\usepackage{amsmath,amssymb}
\setcounter{secnumdepth}{-\maxdimen} % remove section numbering
\usepackage{iftex}
\ifPDFTeX
  \usepackage[T1]{fontenc}
  \usepackage[utf8]{inputenc}
  \usepackage{textcomp} % provide euro and other symbols
\else % if luatex or xetex
  \usepackage{unicode-math} % this also loads fontspec
  \defaultfontfeatures{Scale=MatchLowercase}
  \defaultfontfeatures[\rmfamily]{Ligatures=TeX,Scale=1}
\fi
\usepackage{lmodern}
\ifPDFTeX\else
  % xetex/luatex font selection
\fi
% Use upquote if available, for straight quotes in verbatim environments
\IfFileExists{upquote.sty}{\usepackage{upquote}}{}
\IfFileExists{microtype.sty}{% use microtype if available
  \usepackage[]{microtype}
  \UseMicrotypeSet[protrusion]{basicmath} % disable protrusion for tt fonts
}{}
\usepackage{setspace}
\makeatletter
\@ifundefined{KOMAClassName}{% if non-KOMA class
  \IfFileExists{parskip.sty}{%
    \usepackage{parskip}
  }{% else
    \setlength{\parindent}{0pt}
    \setlength{\parskip}{6pt plus 2pt minus 1pt}}
}{% if KOMA class
  \KOMAoptions{parskip=half}}
\makeatother
% Make \paragraph and \subparagraph free-standing
\makeatletter
\ifx\paragraph\undefined\else
  \let\oldparagraph\paragraph
  \renewcommand{\paragraph}{
    \@ifstar
      \xxxParagraphStar
      \xxxParagraphNoStar
  }
  \newcommand{\xxxParagraphStar}[1]{\oldparagraph*{#1}\mbox{}}
  \newcommand{\xxxParagraphNoStar}[1]{\oldparagraph{#1}\mbox{}}
\fi
\ifx\subparagraph\undefined\else
  \let\oldsubparagraph\subparagraph
  \renewcommand{\subparagraph}{
    \@ifstar
      \xxxSubParagraphStar
      \xxxSubParagraphNoStar
  }
  \newcommand{\xxxSubParagraphStar}[1]{\oldsubparagraph*{#1}\mbox{}}
  \newcommand{\xxxSubParagraphNoStar}[1]{\oldsubparagraph{#1}\mbox{}}
\fi
\makeatother

\usepackage{color}
\usepackage{fancyvrb}
\newcommand{\VerbBar}{|}
\newcommand{\VERB}{\Verb[commandchars=\\\{\}]}
\DefineVerbatimEnvironment{Highlighting}{Verbatim}{commandchars=\\\{\}}
% Add ',fontsize=\small' for more characters per line
\usepackage{framed}
\definecolor{shadecolor}{RGB}{241,243,245}
\newenvironment{Shaded}{\begin{snugshade}}{\end{snugshade}}
\newcommand{\AlertTok}[1]{\textcolor[rgb]{0.68,0.00,0.00}{#1}}
\newcommand{\AnnotationTok}[1]{\textcolor[rgb]{0.37,0.37,0.37}{#1}}
\newcommand{\AttributeTok}[1]{\textcolor[rgb]{0.40,0.45,0.13}{#1}}
\newcommand{\BaseNTok}[1]{\textcolor[rgb]{0.68,0.00,0.00}{#1}}
\newcommand{\BuiltInTok}[1]{\textcolor[rgb]{0.00,0.23,0.31}{#1}}
\newcommand{\CharTok}[1]{\textcolor[rgb]{0.13,0.47,0.30}{#1}}
\newcommand{\CommentTok}[1]{\textcolor[rgb]{0.37,0.37,0.37}{#1}}
\newcommand{\CommentVarTok}[1]{\textcolor[rgb]{0.37,0.37,0.37}{\textit{#1}}}
\newcommand{\ConstantTok}[1]{\textcolor[rgb]{0.56,0.35,0.01}{#1}}
\newcommand{\ControlFlowTok}[1]{\textcolor[rgb]{0.00,0.23,0.31}{\textbf{#1}}}
\newcommand{\DataTypeTok}[1]{\textcolor[rgb]{0.68,0.00,0.00}{#1}}
\newcommand{\DecValTok}[1]{\textcolor[rgb]{0.68,0.00,0.00}{#1}}
\newcommand{\DocumentationTok}[1]{\textcolor[rgb]{0.37,0.37,0.37}{\textit{#1}}}
\newcommand{\ErrorTok}[1]{\textcolor[rgb]{0.68,0.00,0.00}{#1}}
\newcommand{\ExtensionTok}[1]{\textcolor[rgb]{0.00,0.23,0.31}{#1}}
\newcommand{\FloatTok}[1]{\textcolor[rgb]{0.68,0.00,0.00}{#1}}
\newcommand{\FunctionTok}[1]{\textcolor[rgb]{0.28,0.35,0.67}{#1}}
\newcommand{\ImportTok}[1]{\textcolor[rgb]{0.00,0.46,0.62}{#1}}
\newcommand{\InformationTok}[1]{\textcolor[rgb]{0.37,0.37,0.37}{#1}}
\newcommand{\KeywordTok}[1]{\textcolor[rgb]{0.00,0.23,0.31}{\textbf{#1}}}
\newcommand{\NormalTok}[1]{\textcolor[rgb]{0.00,0.23,0.31}{#1}}
\newcommand{\OperatorTok}[1]{\textcolor[rgb]{0.37,0.37,0.37}{#1}}
\newcommand{\OtherTok}[1]{\textcolor[rgb]{0.00,0.23,0.31}{#1}}
\newcommand{\PreprocessorTok}[1]{\textcolor[rgb]{0.68,0.00,0.00}{#1}}
\newcommand{\RegionMarkerTok}[1]{\textcolor[rgb]{0.00,0.23,0.31}{#1}}
\newcommand{\SpecialCharTok}[1]{\textcolor[rgb]{0.37,0.37,0.37}{#1}}
\newcommand{\SpecialStringTok}[1]{\textcolor[rgb]{0.13,0.47,0.30}{#1}}
\newcommand{\StringTok}[1]{\textcolor[rgb]{0.13,0.47,0.30}{#1}}
\newcommand{\VariableTok}[1]{\textcolor[rgb]{0.07,0.07,0.07}{#1}}
\newcommand{\VerbatimStringTok}[1]{\textcolor[rgb]{0.13,0.47,0.30}{#1}}
\newcommand{\WarningTok}[1]{\textcolor[rgb]{0.37,0.37,0.37}{\textit{#1}}}

\usepackage{longtable,booktabs,array}
\usepackage{calc} % for calculating minipage widths
% Correct order of tables after \paragraph or \subparagraph
\usepackage{etoolbox}
\makeatletter
\patchcmd\longtable{\par}{\if@noskipsec\mbox{}\fi\par}{}{}
\makeatother
% Allow footnotes in longtable head/foot
\IfFileExists{footnotehyper.sty}{\usepackage{footnotehyper}}{\usepackage{footnote}}
\makesavenoteenv{longtable}
\usepackage{graphicx}
\makeatletter
\newsavebox\pandoc@box
\newcommand*\pandocbounded[1]{% scales image to fit in text height/width
  \sbox\pandoc@box{#1}%
  \Gscale@div\@tempa{\textheight}{\dimexpr\ht\pandoc@box+\dp\pandoc@box\relax}%
  \Gscale@div\@tempb{\linewidth}{\wd\pandoc@box}%
  \ifdim\@tempb\p@<\@tempa\p@\let\@tempa\@tempb\fi% select the smaller of both
  \ifdim\@tempa\p@<\p@\scalebox{\@tempa}{\usebox\pandoc@box}%
  \else\usebox{\pandoc@box}%
  \fi%
}
% Set default figure placement to htbp
\def\fps@figure{htbp}
\makeatother





\setlength{\emergencystretch}{3em} % prevent overfull lines

\providecommand{\tightlist}{%
  \setlength{\itemsep}{0pt}\setlength{\parskip}{0pt}}



 


\usepackage{booktabs}
\usepackage{longtable}
\usepackage{array}
\usepackage{multirow}
\usepackage{wrapfig}
\usepackage{float}
\usepackage{colortbl}
\usepackage{pdflscape}
\usepackage{tabu}
\usepackage{threeparttable}
\usepackage{threeparttablex}
\usepackage[normalem]{ulem}
\usepackage{makecell}
\usepackage{xcolor}
\makeatletter
\@ifpackageloaded{caption}{}{\usepackage{caption}}
\AtBeginDocument{%
\ifdefined\contentsname
  \renewcommand*\contentsname{Table of contents}
\else
  \newcommand\contentsname{Table of contents}
\fi
\ifdefined\listfigurename
  \renewcommand*\listfigurename{List of Figures}
\else
  \newcommand\listfigurename{List of Figures}
\fi
\ifdefined\listtablename
  \renewcommand*\listtablename{List of Tables}
\else
  \newcommand\listtablename{List of Tables}
\fi
\ifdefined\figurename
  \renewcommand*\figurename{Figure}
\else
  \newcommand\figurename{Figure}
\fi
\ifdefined\tablename
  \renewcommand*\tablename{Table}
\else
  \newcommand\tablename{Table}
\fi
}
\@ifpackageloaded{float}{}{\usepackage{float}}
\floatstyle{ruled}
\@ifundefined{c@chapter}{\newfloat{codelisting}{h}{lop}}{\newfloat{codelisting}{h}{lop}[chapter]}
\floatname{codelisting}{Listing}
\newcommand*\listoflistings{\listof{codelisting}{List of Listings}}
\makeatother
\makeatletter
\makeatother
\makeatletter
\@ifpackageloaded{caption}{}{\usepackage{caption}}
\@ifpackageloaded{subcaption}{}{\usepackage{subcaption}}
\makeatother
\usepackage{bookmark}
\IfFileExists{xurl.sty}{\usepackage{xurl}}{} % add URL line breaks if available
\urlstyle{same}
\hypersetup{
  pdftitle={Supplemental Information},
  colorlinks=true,
  linkcolor={blue},
  filecolor={Maroon},
  citecolor={Blue},
  urlcolor={Blue},
  pdfcreator={LaTeX via pandoc}}


\title{Supplemental Information}
\author{}
\date{}
\begin{document}
\maketitle


\setstretch{1.05}
\section{Supplemental Methods}\label{supplemental-methods}

\subsection{Mixed effects model
details}\label{mixed-effects-model-details}

\subsection{Efficiency computation}\label{efficiency-computation}

\section{Supplemental Results}\label{supplemental-results}

\begin{figure}

\centering{

\includegraphics[width=1\linewidth,height=\textheight,keepaspectratio]{supplemental_files/figure-pdf/fig-cross-sectional-confound-1.pdf}

}

\caption{\label{fig-cross-sectional-confound}Non-identifiability of
growth models in cross-sectional data. Panels show the same simulated
data, reflecting the possibility of stable between-child intercept
differences (A) vs.~varying growth slopes (B).}

\end{figure}%

\begin{figure}

\centering{

\includegraphics[width=1\linewidth,height=\textheight,keepaspectratio]{supplemental_files/figure-pdf/fig-glmer-ladder-full-1.pdf}

}

\caption{\label{fig-glmer-ladder-full}}

\end{figure}%

\begin{Shaded}
\begin{Highlighting}[]
\FunctionTok{\#\# Relation to prior accumulator models  [draft for promotion — MCF to review]}

\NormalTok{Our model is one member of a family of accumulator models of early}
\NormalTok{vocabulary, and the family is easiest to compare in a single set of}
\NormalTok{coordinates. Writing the latent ability of child $i$ on word $j$ at}
\NormalTok{age $t$ as}

\NormalTok{$$\textbackslash{}eta\_\{ij\}(t) = \textbackslash{}log r\_i + \textbackslash{}log \textbackslash{}alpha\_i + \textbackslash{}kappa\_i \textbackslash{}log t {-} \textbackslash{}delta\_j,$$}

\NormalTok{every prior model in this tradition is a special case that fixes one or}
\NormalTok{more of three quantities: the between{-}child variance in efficiency}
\NormalTok{($\textbackslash{}sigma\_\textbackslash{}alpha$), the population scaling exponent ($\textbackslash{}kappa\_\{\textbackslash{}text\{pop\}\}$),}
\NormalTok{and — critically — the between{-}child variance in that exponent}
\NormalTok{($\textbackslash{}sigma\_\textbackslash{}kappa$). Mapping each prior model into these coordinates makes}
\NormalTok{the comparison structural rather than notational.}

\NormalTok{**McMurray [{-}@mcmurray2007]** is a pure accumulator: every word fills at a}
\NormalTok{single unit rate ($\textbackslash{}kappa = 1$) and children do not differ}
\NormalTok{($\textbackslash{}sigma\_\textbackslash{}alpha = \textbackslash{}sigma\_\textbackslash{}kappa = 0$). Its central insight is that the}
\NormalTok{*vocabulary* curve accelerates even when the underlying rate does not,}
\NormalTok{because ability sweeps through a distribution of word difficulties — the}
\NormalTok{construction reproduced in panel (A) of our schematic. Acceleration here}
\NormalTok{is a population{-}level property of the difficulty distribution, not a}
\NormalTok{change in the child. This corresponds to our M1.}

\NormalTok{**Mitchell \& McMurray [{-}@mitchell2009]** add a leverage term: already{-}known}
\NormalTok{words speed the acquisition of new ones. They prove that this feedback}
\NormalTok{shifts *when* words are acquired but cannot manufacture acceleration that}
\NormalTok{is not already present in the difficulty distribution. In our coordinates}
\NormalTok{it moves the timing of acquisition without changing $\textbackslash{}kappa$, and — being}
\NormalTok{non{-}linear in accumulated vocabulary — it is the one prior model that does}
\NormalTok{not embed cleanly in the linear predictor above.}

\NormalTok{**Hidaka [{-}@hidaka2013]** introduces a rate{-}change process in which ability}
\NormalTok{grows as $t\^{}\{D+1\}$; his exponent $D$ is exactly $\textbackslash{}kappa\_\{\textbackslash{}text\{pop\}\} {-} 1$}
\NormalTok{in our notation, so $D \textgreater{} 0$ is population{-}level acceleration}
\NormalTok{($\textbackslash{}kappa\_\{\textbackslash{}text\{pop\}\} \textgreater{} 1$). This is our M2. The acceleration is shared}
\NormalTok{across children ($\textbackslash{}sigma\_\textbackslash{}kappa = 0$). Hidaka relates $D$ to word}
\NormalTok{imageability — i.e. a *per{-}word* scaling exponent — which is an}
\NormalTok{interesting direction we do not pursue here.}

\NormalTok{**Mollica \& Piantadosi [{-}@mollica2017]** model each word as a Gamma}
\NormalTok{waiting{-}time process in which acquisition occurs after roughly ten}
\NormalTok{"effective learning instances." In our coordinates this is again a pure}
\NormalTok{accumulator ($\textbackslash{}kappa = 1$); the per{-}word variation they capture lives in}
\NormalTok{the Gamma shape parameter, which corresponds to a per{-}word discrimination}
\NormalTok{(a 2PL extension) rather than to per{-}child structure.}

\NormalTok{**Kachergis et al. [{-}@kachergis2022]** are the first to free}
\NormalTok{$\textbackslash{}sigma\_\textbackslash{}alpha$: they cast vocabulary as a Rasch model in which each child}
\NormalTok{has their own latent ability, unifying the accumulator and item{-}response}
\NormalTok{traditions. This is our M3. But the scaling exponent remains shared across}
\NormalTok{children ($\textbackslash{}sigma\_\textbackslash{}kappa = 0$), and the model is proposed as a framework}
\NormalTok{rather than fit to longitudinal individual data.}

\NormalTok{Reading down this list, the literature progressively frees parameters —}
\NormalTok{McMurray fixes all three, Hidaka frees $\textbackslash{}kappa\_\{\textbackslash{}text\{pop\}\}$, Kachergis}
\NormalTok{frees $\textbackslash{}sigma\_\textbackslash{}alpha$ — but every model holds $\textbackslash{}sigma\_\textbackslash{}kappa = 0$.}

\PreprocessorTok{|}\NormalTok{ Model }\PreprocessorTok{|}\NormalTok{ $\textbackslash{}sigma\_\textbackslash{}alpha$ }\PreprocessorTok{|}\NormalTok{ $\textbackslash{}kappa\_\{\textbackslash{}text\{pop\}\}$ }\PreprocessorTok{|}\NormalTok{ $\textbackslash{}sigma\_\textbackslash{}kappa$ }\PreprocessorTok{|}\NormalTok{ per{-}word $\textbackslash{}delta\_j$ }\PreprocessorTok{|}
\PreprocessorTok{|{-}{-}{-}|}\NormalTok{:{-}{-}:}\PreprocessorTok{|}\NormalTok{:{-}{-}:}\PreprocessorTok{|}\NormalTok{:{-}{-}:}\PreprocessorTok{|}\NormalTok{:{-}{-}:}\PreprocessorTok{|}
\PreprocessorTok{|}\NormalTok{ McMurray (2007)            }\PreprocessorTok{|}\NormalTok{ 0 }\PreprocessorTok{|}\NormalTok{ 1     }\PreprocessorTok{|}\NormalTok{ 0 }\PreprocessorTok{|}\NormalTok{ free }\PreprocessorTok{|}
\PreprocessorTok{|}\NormalTok{ Mitchell \& McMurray (2009) }\PreprocessorTok{|}\NormalTok{ 0 }\PreprocessorTok{|}\NormalTok{ 1     }\PreprocessorTok{|}\NormalTok{ 0 }\PreprocessorTok{|}\NormalTok{ free + leverage }\PreprocessorTok{|}
\PreprocessorTok{|}\NormalTok{ Hidaka (2013), rate{-}change }\PreprocessorTok{|}\NormalTok{ 0 }\PreprocessorTok{|}\NormalTok{ $D+1$ }\PreprocessorTok{|}\NormalTok{ 0 }\PreprocessorTok{|}\NormalTok{ free }\PreprocessorTok{|}
\PreprocessorTok{|}\NormalTok{ Mollica \& Piantadosi (2017)}\PreprocessorTok{|}\NormalTok{ 0 }\PreprocessorTok{|}\NormalTok{ 1     }\PreprocessorTok{|}\NormalTok{ 0 }\PreprocessorTok{|}\NormalTok{ free }\PreprocessorTok{|}
\PreprocessorTok{|}\NormalTok{ Kachergis et al. (2022)    }\PreprocessorTok{|}\NormalTok{ free }\PreprocessorTok{|}\NormalTok{ free }\PreprocessorTok{|}\NormalTok{ 0 }\PreprocessorTok{|}\NormalTok{ free }\PreprocessorTok{|}
\PreprocessorTok{|}\NormalTok{ **This work (M4)**         | **free** | **free** | **free** | **free** }\PreprocessorTok{|}

\NormalTok{Our contribution is to free $\textbackslash{}sigma\_\textbackslash{}kappa$: children differ not only in}
\NormalTok{how efficient they are but in how steeply their learning accelerates. This}
\NormalTok{is the load{-}bearing addition — it is what lets the model fit the widening}
\NormalTok{fan of trajectories with age (M4 vs M3 in the main text). It is also a}
\NormalTok{falsifiable claim rather than a parameter convenience: if $\textbackslash{}sigma\_\textbackslash{}kappa$}
\NormalTok{were near zero the model would collapse onto Kachergis et al., but the}
\NormalTok{estimated per{-}child slope variance is well above zero and large relative}
\NormalTok{to the efficiency variance $\textbackslash{}sigma\_\textbackslash{}alpha$ (main text, Figure 2 and}
\NormalTok{Results).}
\end{Highlighting}
\end{Shaded}

\begin{Shaded}
\begin{Highlighting}[]
\FunctionTok{\#\# Comparing children and language models on one slope axis  [draft for promotion — MCF to review]}

\NormalTok{The comparison to language models in Figure 6 places children and models}
\NormalTok{on a single, literally shared axis: the slope of the acquisition sigmoid}
\NormalTok{in log{-}experience. Both populations have sigmoidal per{-}word acquisition}
\NormalTok{curves, and the slope of that sigmoid is directly comparable.}

\NormalTok{On the child side, our linear predictor makes this slope exact.}
\NormalTok{Differentiating $\textbackslash{}eta\_\{ij\}(t)$ with respect to $\textbackslash{}log t$ gives $\textbackslash{}kappa\_i$:}
\NormalTok{the logit{-}probability that a child produces a word rises linearly in}
\NormalTok{log{-}age with slope equal to that child\textquotesingle{}s scaling exponent. The per{-}child}
\NormalTok{sigmoid slope *is* $\textbackslash{}kappa\_i$, by construction.}

\NormalTok{On the model side, @chang2022 fit a four{-}parameter logistic to each}
\NormalTok{(model, word) pair, describing how the model\textquotesingle{}s surprisal for a word falls}
\NormalTok{over training. Re{-}expressed as the fraction of a word\textquotesingle{}s surprisal}
\NormalTok{reduction achieved so far, their fit is a sigmoid in}
\NormalTok{$\textbackslash{}log\_\{10\}(\textbackslash{}text\{training steps\})$ with slope $1/\textbackslash{}text\{scale\}\_w$.}
\NormalTok{Dividing by $\textbackslash{}ln 10$ converts this to natural{-}log units — the same units}
\NormalTok{as $\textbackslash{}kappa\_i$ — so the two slopes can be read off the same axis.}

\NormalTok{The result is a striking separation (Figure 6). Language models cluster}
\NormalTok{near a slope of one — the unit{-}accumulator boundary, where each log{-}unit}
\NormalTok{of additional experience yields a fixed increment of learning. Children}
\NormalTok{sit in a broad distribution roughly an order of magnitude higher. On the}
\NormalTok{shared axis, children\textquotesingle{}s word learning accelerates with experience in a way}
\NormalTok{that these models\textquotesingle{} does not.}

\NormalTok{Three caveats bound this comparison. First, the variability is asymmetric:}
\NormalTok{children\textquotesingle{}s slopes vary across *children*, whereas a single model is a}
\NormalTok{deterministic learner, so model slopes vary across *words* within a run}
\NormalTok{rather than across replicate learners. Second, the steep child slope is}
\NormalTok{the *per{-}child* quantity; a researcher fitting a single sigmoid to}
\NormalTok{cross{-}sectional Wordbank data per word would recover a shallower slope,}
\NormalTok{because marginalizing over the population attenuates it. Third, the models}
\NormalTok{in @chang2022 were trained on adult written text rather than}
\NormalTok{child{-}directed speech, so part of the gap could reflect the input}
\NormalTok{distribution rather than the learner. Models trained directly on CHILDES}
\NormalTok{(Feng et al., forthcoming) provide the matched{-}input test: if the gap}
\NormalTok{largely persists, it is structural; if it shrinks substantially, the input}
\NormalTok{distribution is doing more of the work than we currently attribute to it.}

\CommentTok{\textless{}!{-}{-} [BIB NEEDED] Feng et al. (CHILDES{-}trained GPT) — no verified entry yet;}
\CommentTok{     named in text only. Daunizeau (2017) logistic{-}Gaussian approximation is}
\CommentTok{     used in the derivation note but not needed for this prose. {-}{-}\textgreater{}}
\end{Highlighting}
\end{Shaded}

\begin{Shaded}
\begin{Highlighting}[]
\FunctionTok{\#\# Connection to neural scaling laws  [draft for promotion — MCF to review; most speculative of the three]}

\NormalTok{The per{-}word picture above connects to a more familiar aggregate}
\NormalTok{phenomenon: neural scaling laws. A model\textquotesingle{}s loss falls as a smooth power}
\NormalTok{law in the amount of training data, and this descent is remarkably}
\NormalTok{regular — different runs, and different model sizes, trace nearly the same}
\NormalTok{curve (Kaplan et al., 2020; Hoffmann et al., 2022). Our per{-}word analysis}
\NormalTok{is the micro{-}level counterpart of that regularity. Each word is acquired}
\NormalTok{along a sigmoid whose slope sits near the unit{-}accumulator value, and the}
\NormalTok{spread of those slopes — across words and across runs — is small. The}
\NormalTok{aggregate smoothness of scaling laws and the per{-}word $\textbackslash{}kappa \textbackslash{}approx 1$}
\NormalTok{we recover are two views of the same property: language models convert}
\NormalTok{experience into knowledge as near{-}pure accumulators, with little}
\NormalTok{instance{-}to{-}instance variation.}

\NormalTok{Children depart from this picture on both axes at once. Their per{-}instance}
\NormalTok{scaling is far steeper ($\textbackslash{}kappa \textbackslash{}approx 10$), and they vary substantially}
\NormalTok{from one another in that exponent ($\textbackslash{}sigma\_\textbackslash{}kappa$ well above zero). The}
\NormalTok{"data gap" between children and language models }\CommentTok{[}\OtherTok{@frank2023; @warstadt2023}\CommentTok{]}
\NormalTok{is therefore not only that children reach adult competence from orders of}
\NormalTok{magnitude less input. It is that children convert experience into knowledge}
\NormalTok{through a different scaling relationship — steeper, and far more variable}
\NormalTok{across individuals — than the one that scaling laws describe for models.}

\NormalTok{This connection is analogical rather than formal: scaling laws characterize}
\NormalTok{aggregate loss, whereas our $\textbackslash{}kappa$ characterizes per{-}item acquisition,}
\NormalTok{and the two need not move together. We offer it as a framing for why the}
\NormalTok{data gap may be harder to close than raw data counts suggest.}

\CommentTok{\textless{}!{-}{-} [BIB NEEDED] Kaplan et al. (2020) and Hoffmann et al. (2022) are named}
\CommentTok{     in text but have NO verified bib entries yet — do not promote with live}
\CommentTok{     @ markers until added. {-}{-}\textgreater{}}
\end{Highlighting}
\end{Shaded}

\begin{figure}

\centering{

\includegraphics[width=1\linewidth,height=\textheight,keepaspectratio]{supplemental_files/figure-pdf/fig-schematic-theta-1.pdf}

}

\caption{\label{fig-schematic-theta}}

\end{figure}%

\begingroup\fontsize{8}{10}\selectfont

\begin{longtable}[t]{llrrrrr}

\caption{\label{tbl-bic}}

\tabularnewline

\toprule
\multicolumn{2}{c}{ } & \multicolumn{5}{c}{$\Delta$BIC} \\
\cmidrule(l{3pt}r{3pt}){3-7}
Model & Age & Thal & Smith & Marchman & Norwegian & Japanese\\
\midrule
M1. accumulator & — & 703,172 & 273,687 & 251,166 & 2,281,553 & 94,461\\
M2. + acceleration & linear & 206,517 & 179,838 & 134,563 & 1,377,575 & 41,138\\
M2. + acceleration & log & 204,895 & 178,437 & 133,843 & 1,347,973 & 39,023\\
M3. + efficiency var. & linear & 35,832 & 17,396 & 13,729 & 130,599 & 3,829\\
M3. + efficiency var. & log & 34,195 & 15,333 & 13,113 & 115,089 & 3,206\\
\addlinespace
M4. + acceleration var. & linear & 1,134 & 553 & 0 & 4,301 & 169\\
M4. + acceleration var. & log & 0 & 0 & 30 & 0 & 0\\
\bottomrule

\end{longtable}

\endgroup{}

\begingroup\fontsize{9}{11}\selectfont

\begin{longtable}[t]{lrrr}

\caption{\label{tbl-blup-dip}}

\tabularnewline

\toprule
Dataset & N & Efficiency (\$\textbackslash{}xi\$) & Acceleration (\$\textbackslash{}zeta\$)\\
\midrule
English (Thal) & 651 & 0.96 & >0.99\\
English (Smith) & 312 & 0.30 & 0.99\\
English (Marchman) & 313 & 0.99 & 0.99\\
Norwegian & 1667 & >0.99 & 0.86\\
Japanese & 176 & 0.85 & 0.99\\
\addlinespace
Pooled & 3119 & 0.98 & 0.96\\
\bottomrule

\end{longtable}

\endgroup{}

\begin{figure}

\centering{

\includegraphics[width=1\linewidth,height=\textheight,keepaspectratio]{supplemental_files/figure-pdf/fig-blup-hist-1.pdf}

}

\caption{\label{fig-blup-hist}}

\end{figure}%

\begingroup\fontsize{8}{10}\selectfont

\begin{longtable}[t]{lllrrrrrr}

\caption{\label{tbl-input-rates}}

\tabularnewline

\toprule
Source & Sample & Channel & n & tok/hr & \$\textbackslash{}pm\$1 SD & tok/mo & \$\textbackslash{}sigma\_r\$ & 1-SD\\
\midrule
Hart \textbackslash{}\& Risley & all kids pooled & mother CDS & 42 & 1,148 & 546--2,410 & 419K & 0.74 & 2.10\$\textbackslash{}times\$\\
Hart \textbackslash{}\& Risley & within-stratum pooled & mother CDS & 42 & 1,148 & 631--2,086 & 419K & 0.60 & 1.82\$\textbackslash{}times\$\\
Soderstrom \textbackslash{}\& Wittebolle & all kids pooled & all speech & 12 & 967 & 622--1,502 & 353K & 0.44 & 1.55\$\textbackslash{}times\$\\
Sperry & all kids pooled & adult-to-child & 42 & 1,537 & 901--2,622 & 561K & 0.53 & 1.71\$\textbackslash{}times\$\\
Sperry & within-stratum pooled & adult-to-child & 42 & 1,537 & 974--2,425 & 561K & 0.46 & 1.58\$\textbackslash{}times\$\\
\addlinespace
Weisleder \textbackslash{}\& Fernald & all kids pooled & mother CDS & 29 & 458 & 197--1,065 & 167K & 0.84 & 2.33\$\textbackslash{}times\$\\
Bergelson 2018 (SEEDLingS, LENA) & US daylong, 6--7 mo & all adult (AWC) & 44 & 1,101 & 767--1,580 & 402K & 0.36 & 1.43\$\textbackslash{}times\$\\
Casillas 2020 (Tseltal Mayan) & Chiapas daylong & child-directed & 10 & 720 & — & 263K & — & —\\
Coffey et al. 2024 (plausible envelope) & cross-cultural bounds & CDS to all-input & — & — & 4--7,232 & — & — & —\\
\bottomrule

\end{longtable}

\endgroup{}

\begin{Shaded}
\begin{Highlighting}[]
\FunctionTok{\#\# Estimating the input{-}rate prior: units and conversions  [draft for promotion — MCF to review]}

\NormalTok{The model represents each child\textquotesingle{}s input as a log rate,}
\NormalTok{$\textbackslash{}log r\_i \textbackslash{}sim \textbackslash{}mathcal\{N\}(\textbackslash{}mu\_r, \textbackslash{}sigma\_r\^{}2)$, measured in **adult tokens}
\NormalTok{per waking hour**. @tbl{-}input{-}rates collects published estimates of this}
\NormalTok{rate, expressed in shared units so they can be compared and so that the}
\NormalTok{implied between{-}child spread $\textbackslash{}sigma\_r$ is visible on one scale.}

\NormalTok{Conventions used throughout:}

\SpecialStringTok{{-} }\NormalTok{**Tokens/hour to tokens/month.** We multiply by $H = 365$ waking hours}
\NormalTok{  per month ($12\textbackslash{}text\{ hr/day\} \textbackslash{}times 30.44\textbackslash{}text\{ days/mo\}$), the model\textquotesingle{}s}
\NormalTok{  constant.}
\SpecialStringTok{{-} }\NormalTok{**$\textbackslash{}sigma\_r$ is on the log scale** — the SD of $\textbackslash{}log(\textbackslash{}text\{tokens/hr\})$}
\NormalTok{  across children within a sample. Its exponential $e\^{}\{\textbackslash{}sigma\_r\}$ is the}
\NormalTok{  multiplicative factor for a one{-}SD shift in input (e.g.}
\NormalTok{  $\textbackslash{}sigma\_r = 0.53 \textbackslash{}Rightarrow 1.7\textbackslash{}times$), and the 16th–84th percentile}
\NormalTok{  spread is $e\^{}\{2\textbackslash{}sigma\_r\}$.}
\SpecialStringTok{{-} }\NormalTok{**Channel.** Studies count different speaker sets; we take the channel}
\NormalTok{  closest to "tokens reaching the child" — adult{-}to{-}child for Sperry,}
\NormalTok{  all{-}adult LENA word counts for SEEDLingS, and maternal CDS for Hart \&}
\NormalTok{  Risley and Weisleder \& Fernald.}
\SpecialStringTok{{-} }\NormalTok{**Pooled vs. within{-}stratum.** For Hart \& Risley and Sperry we report}
\NormalTok{  $\textbackslash{}sigma\_r$ both across all children (which mixes in SES and site}
\NormalTok{  gradients) and within stratum (subtracting each band\textquotesingle{}s mean first); the}
\NormalTok{  latter is the cleaner within{-}population estimate.}
\SpecialStringTok{{-} }\NormalTok{**Minutes to tokens.** Where a source reports minutes of speech per hour}
\NormalTok{  rather than token counts (Casillas; the Coffey bounds), we convert at}
\NormalTok{  $\textbackslash{}approx 200$ words/min ($0.3$ s/word, the Brent{-}corpus CDS rate used by}
\NormalTok{  @coffey2024); those rows are accurate to roughly $\textbackslash{}pm 20\textbackslash{}\%$.}

\NormalTok{Rows marked *computed* are recomputed here from per{-}dyad data; *literature*}
\NormalTok{rows are taken from the curated validation set in }\InformationTok{\textasciigrave{}input\_estimation/\textasciigrave{}}\NormalTok{.}
\end{Highlighting}
\end{Shaded}

\begin{Shaded}
\begin{Highlighting}[]
\FunctionTok{\#\# Analytic input{-}rate ($\textbackslash{}sigma\_r$) sensitivity  [draft for promotion — MCF to review]}

\NormalTok{The input/efficiency variance split is identified only by the external pin on}
\NormalTok{the population SD of log input rate, $\textbackslash{}sigma\_r$ — we do not observe input in the}
\NormalTok{longitudinal data. Because every reported input share therefore rests on a}
\NormalTok{quantity we assume rather than estimate, we map how the conclusion moves as}
\NormalTok{$\textbackslash{}sigma\_r$ ranges over its plausible values (panel B of @fig{-}io{-}partition).}

\NormalTok{The data identify the *total* between{-}child intercept variance,}
\NormalTok{$\textbackslash{}sigma\_\textbackslash{}xi\^{}2 = \textbackslash{}sigma\_\textbackslash{}alpha\^{}2 + \textbackslash{}sigma\_r\^{}2$ (the spread of children\textquotesingle{}s ability}
\NormalTok{at the reference age $a\_0$), very precisely; the pin only decides how that fixed}
\NormalTok{total is *split* into an efficiency part $\textbackslash{}sigma\_\textbackslash{}alpha\^{}2$ and an input part}
\NormalTok{$\textbackslash{}sigma\_r\^{}2$. Holding $\textbackslash{}sigma\_\textbackslash{}xi\^{}2$ at its fitted value, the input share is then}
\NormalTok{a closed form,}

\NormalTok{$$1{-}\textbackslash{}pi\_\textbackslash{}alpha(\textbackslash{}sigma\_r) = \textbackslash{}frac\{\textbackslash{}sigma\_r\^{}2\}\{\textbackslash{}sigma\_\textbackslash{}xi\^{}2\},$$}

\NormalTok{which we evaluate across the plausible $\textbackslash{}sigma\_r$ range *without refitting* — the}
\NormalTok{analytic sensitivity curve.}

\NormalTok{This is first{-}order: it treats $\textbackslash{}sigma\_\textbackslash{}xi\^{}2$ as itself invariant to the pin.}
\NormalTok{Refits at off{-}pin values reveal a small second{-}order feedback — pinning a larger}
\NormalTok{$\textbackslash{}sigma\_r$ slightly *shrinks* the estimated total variance (for English,}
\NormalTok{$\textbackslash{}sigma\_\textbackslash{}xi\^{}2$ falls from 3.57 at $\textbackslash{}sigma\_r=0.53$ to 3.22 at $\textbackslash{}sigma\_r=0.80$). The}
\NormalTok{analytic curve therefore slightly *understates* the input share at large}
\NormalTok{$\textbackslash{}sigma\_r$ (≈18\% analytic vs. ≈20\% by refit at $\textbackslash{}sigma\_r=0.80$), an error of about}
\NormalTok{two points of $\textbackslash{}pi\_\textbackslash{}alpha$ in the conservative direction for the size of the}
\NormalTok{input effect. The open points in panel B are these refit anchors; the gap}
\NormalTok{between point and curve is the second{-}order term.}

\NormalTok{We use the analytic curve, anchored by the few refits we have, rather than a}
\NormalTok{dense refit sweep: the qualitative conclusion — input accounts for a modest share}
\NormalTok{even at the top of the plausible $\textbackslash{}sigma\_r$ range — is unaffected by the}
\NormalTok{\textasciitilde{}2{-}point correction, whereas a full sweep of the Norwegian model alone would cost}
\NormalTok{roughly half a day of compute per pin value.}
\end{Highlighting}
\end{Shaded}





\end{document}
